% part: intuitionistic-logic
% chapter: introduction
% section: syntax

\documentclass[../../../include/open-logic-section]{subfiles}

\begin{document}

\olfileid{int}{int}{syn}

\olsection{وجدانیت پسند منطق دی نحو}

وجدانیت پسند منطق دی نحو قضیاتی منطق والی ای اے۔ کلاسیکی
قضیاتی منطق وچ کجھ رابطے ہور رابطیاں نال متعین ہو سکدے نیں۔
مثلاً $!A \lif !B$ نوں $\lnot !A \lor !B$ نال، یا $!A \lor !B$
نوں $\lnot(\lnot !A \land \lnot !B)$ نال متعین کر سکدے آں۔
ایس لئی کلاسیکی منطق نوں پیش کرن ویلے کجھ رابطے اکثر ایہناں
تعریفاں دی مختصر لکھت وجوں لیاندے جاندے نیں۔ وجدانیت پسند منطق
وچ ایہ گل دو صورتاں وچ چلدی اے۔ $\lnot !A$ نوں $!A \lif \lfalse$
دی مختصر لکھت وجوں متعین کر سکدے آں تے اکثر کردے آں۔ فیر
$\lfalse$ دی اپنی تعریف نہیں ہونی چاہیدی!{} نالے $!A \liff !B$
نوں کلاسیکی منطق وانگ $(!A \lif !B) \land (!B \lif !A)$
دے طور اُتے متعین کر سکدے آں۔

قضیاتی وجدانیت پسند منطق دے !!^{formula}s،
\emph{!!{propositional variable}s} تے قضیاتی مستقل~$\lfalse$ توں
\emph{منطقی رابطیاں} نال بندے نیں۔ ساڈے کول ایہ نیں:

\begin{enumerate}
\item !!^a{denumerable} سیٹ~$\PVar$، جیہدے !!{propositional variable}s
  $\Obj p_0$، $\Obj p_1$، \dots نیں۔
\item !!{falsity} لئی قضیاتی مستقل~$\lfalse$۔
\item منطقی رابطے $\land$ (عطف)، $\lor$ (فصل)، $\lif$ (شرطیہ)۔
\item رموزِ اوقاف، کھلدا قوس «(»، بند ہون والا قوس «)» تے کوما۔
\end{enumerate}

\begin{defn}[فارمولا]
\ollabel{defn:formulas} قضیاتی وجدانیت پسند منطق دے
\emph{!!{formula}s} دا سیٹ~$\Frm[L_0]$ استقرا نال ایویں متعین ہوندا اے:
\begin{enumerate}
\item $\lfalse$ اک ایٹمی !!{formula} اے۔

\item ہر !!{propositional variable}~$\Obj p_i$ اک ایٹمی
  !!{formula} اے۔

\item جے $!A$ تے $!B$، !!{formula}s نیں، تاں $(!A \land
  !B)$ اک !!a{formula} اے۔

\item جے $!A$ تے $!B$، !!{formula}s نیں، تاں $(!A \lor !B)$
  اک !!a{formula} اے۔

\item جے $!A$ تے $!B$، !!{formula}s نیں، تاں $(!A \lif !B)$
  اک !!a{formula} اے۔

\tagitem{limitClause}{ہور کوئی شے !!a{formula} نہیں۔}{}
\end{enumerate}
\end{defn}

اُتے دتے ابتدائی رابطیاں توں علاوہ، اسیں ایہ \emph{متعین} نشان
وی ورتدے آں، $\lnot$ (نفی) تے $\liff$ (!!{biconditional})۔
ایہناں متعین عاملاں نال بنے فارمولے ایویں سمجھے جان:
\begin{enumerate}
\item $\lnot !A$، $!A \lif \lfalse$ دی مختصر لکھت اے۔

\item $!A \liff !B$، $(!A \lif !B) \land (!B
  \lif !A)$ دی مختصر لکھت اے۔
\end{enumerate}

رسمی طور اُتے $\lnot$ اک مختصر لکھت اے، پر کدی کدی اسیں
~$\lnot$ لئی تعریف وچ صاف شقّاں تے قاعدے دیاں گے، جیویں ایہ
ابتدائی نشان ہووے۔ ایہ زیادہ تر ایس لئی اے کہ اسیں مشق دے
مسئلے بیان کر سکیے۔

\end{document}
